\documentclass[10pt,a4paper]{amsart}
\usepackage[margin=19mm]{geometry}
\usepackage{amsmath,amssymb,mathtools}
\usepackage[scr=rsfs,cal=euler]{mathalfa}
\usepackage{xfrac}
\usepackage[unicode,hidelinks,bookmarks=false]{hyperref}
\newcommand{\ie}{i.e.\ }
\newcommand{\cf}{cf.\ }
\newcommand{\resp}{respectively\ }
\newcommand{\Subsection}{\subsection}
\providecommand{\nobreakdash}{\nobreak\hbox{-}}
\setlength{\emergencystretch}{2em}
\multlinegap0pt
\usepackage{bm}
\usepackage{bbm}
\usepackage{dsfont}
\usepackage{xspace}
\usepackage{cancel}
\usepackage{mathtools}
\usepackage{amsthm}
\makeatletter
\@ifundefined{theorem}{\newtheorem{theorem}{Theorem}}{}
\@ifundefined{proposition}{\newtheorem{proposition}[theorem]{Proposition}}{}
\makeatother
\newcommand{\R}{\mathbb{R}}
\newcommand{\dd}{\mathrm{d}}
\title[Reduction of Multiple Orthogonal Polynomials to Standard Ort]{Reduction of Multiple Orthogonal Polynomials to Standard Orthogonal Polynomials}
\author{Roozbeh Gharakhloo, Maksim Kosmakov, Kenta Miyahara}
\date{Source: arXiv:2606.27594v1 (CC BY 4.0)}
\begin{document}
\maketitle

\noindent\emph{Source and licence.} Reduction of Multiple Orthogonal Polynomials to Standard Orthogonal Polynomials. Version arXiv:2606.27594v1 (CC BY 4.0), distributed under the Creative Commons Attribution 4.0 International licence, as recorded in the arXiv licence metadata for this exact version and shown on the abstract page https://arxiv.org/abs/2606.27594v1. Full rights notice: licenses/arxiv-2606.27594-CC-BY-4.0.txt.\par\smallskip

\noindent\textit{Editorial scope.} Counted unit: the Proposition whose source label is \texttt{app bordered hankel prop} in the article (source0.tex lines 2405--2425, source label \texttt{app bordered hankel prop}), delivered together with the article's own complete original proof of it (source0.tex lines 2427--2488, one \texttt{proof} environment of 62 lines). No mathematics is written, repaired or re-derived: statement and proof are the article's own text line for line. Required context retained uncounted: defn of H (lines 359--361); defn of M, H 4 blocks (lines 387--395); app bordered hankel vector (lines 2392--2394); app bordered hankel measures (lines 2396--2398); eq:orthog (lines 2400--2402). Every definition, display and label of those lines is retained unchanged. Omitted from this package and not redistributed: 1--358, 362--386, 396--2391, 2395--2395, 2399--2399, 2403--2404, 2426--2426, 2489--2647. Nothing inside a retained excerpt is omitted or altered: every retained body is delivered exactly as the article's lines are written, so no omission receipt of any kind accompanies this package. Citations: the retained excerpts carry no citation command at all, so no bibliography entry of the article is transcribed and no reference is redistributed. Reference adaptation: none was needed and none was made; every reference inside the delivered unit resolves inside it, so no unresolved reference or citation is produced. Printed slips kept literal: no printed word, symbol, hypothesis or number of the counted statement or of its proof was altered. Layout and class: the article's own class is \texttt{amsart} with packages including amsmath, amsthm, amssymb, stmaryrd, url, here, xy, comment, mathtools, ascmac, mathrsfs, bm, ulem, framed; the wrapper is the lane's pinned \texttt{amsart} editorial wrapper at 10pt on A4, which restates the theorem-like environments the retained text needs and adds the article's own macro definitions that the retained text needs (\texttt{\textbackslash R}, \texttt{\textbackslash dd}), with extra wrapper packages (bm, bbm, dsfont, xspace, cancel, mathtools). The delivered numbering is the unit's own. Assets: no retained span contains a figure environment, a graphics-inclusion command, a PDF-page inclusion, a table or a TikZ picture; no image file is copied into this package, the package declares no asset dependency and no image file is redistributed with it. Licence: version arXiv:2606.27594v1 of this article is distributed under the Creative Commons Attribution 4.0 International licence, as recorded in the arXiv licence metadata for this exact version and shown on the abstract page https://arxiv.org/abs/2606.27594v1. A full rights notice is delivered at licenses/arxiv-2606.27594-CC-BY-4.0.txt. Characters: every retained excerpt is pure ASCII, so the delivered engine, XeLaTeX with its default Latin Modern text font, reports no missing character.
\begin{equation} \label{defn of H}
    \bm{H} = \operatorname{diag}(h_0,h_1,\ldots, h_{n-1}).
\end{equation}

\begin{equation} \label{defn of M, H 4 blocks}
\bm{M} =
\begin{pmatrix}
\bm{M}_{11} & \bm{M}_{12} \\
\bm{M}_{21} & \bm{M}_{22}
\end{pmatrix},
\qquad
\bm{H} = \operatorname{diag}(\bm{H}_1,\bm{H}_2).
\end{equation}

\begin{align} \label{app bordered hankel vector}
\vec{n} = (n-m, 1, 1, \cdots, 1) \in \R^n,
\end{align}

\begin{align} \label{app bordered hankel measures}
\vec{\mu}(x) = \left( \mu(x), \delta(x - z_1), \delta(x - z_2),  \cdots, \delta(x - z_m) \right),    
\end{align}

\begin{equation}\label{eq:orthog}
\int_{\mathscr I} p_i(x)\, p_j(x)\,\dd\mu(x)=h_i\,\delta_{ij},\quad h_i>0.
\end{equation}

\begin{proposition} \label{app bordered hankel prop}
The determinant of the bordered Hankel matrix
\begin{align} \label{borderd hankel defn}
\bm H_n^{B}[\mu;z_1,\cdots,z_m] = \begin{pmatrix}
		\mu_0 & \mu_1 & \cdots & \mu_{n-m-1} & 1 & 1 & \cdots & 1 \\
		\mu_1 & \mu_2 & \cdots & \mu_{n-m} & z_1 & z_2 & \cdots & z_m \\
		\mu_2 & \mu_3 & \cdots & \mu_{n-m+1} & z_1^2 & z_2^2 & \cdots & z_m^2 \\
		\vdots & \vdots & \reflectbox{$\ddots$} & \vdots & \vdots & \vdots & \ddots & \vdots  \\
		\mu_{n-1} & \mu_n & \cdots & \mu_{2n-m-2} & z_1^{n-1} & z_2^{n-1} & \cdots & z_{m}^{n-1}\end{pmatrix}
\end{align}
satisfies
\[
\det\bm H_n^{B}[\mu;z_1,\cdots,z_m] = \det \bm{H}_{n-m}[\mu] \det 
\begin{pmatrix}
	p_{n-m}(z_1) & \cdots & p_{n-m}(z_{m}) \\
	\vdots & \ddots & \vdots \\
	p_{n-1}(z_1) & \cdots & p_{n-1}(z_{m})
\end{pmatrix}
\]
where \[\det \bm{H}_{n-m}[\mu] = \prod_{k=0}^{n-m-1} h_k.\]
\end{proposition}

\begin{proof}
We define two $n$-column vectors,
\[
\bm p_{n}(x) = \begin{pmatrix}
    \bm p_{1}(x)\\
    \bm p_{2}(x)
\end{pmatrix}, \quad \bm m_{n}(x) = \begin{pmatrix}
    \bm m_{1}(x)\\
    \bm m_{2}(x)
\end{pmatrix},
\]
where
\begin{align}
&\bm p_{1}(x) = \left(p_0(x), p_1(x), \cdots,p_{n-m-1}(x) \right)^\top, \quad \bm p_{2}(x) = \left( p_{n-m}(x), p_{n-m+1}(x), \cdots, p_{n-1}(x) \right)^\top,\\
&\bm m_{1}(x) = \left( 1,x,\cdots,x^{n-m-1} \right)^\top, \quad \bm m_{2}(x) = \left( x^{n-m}, x^{n-m+1}, \cdots, x^{n-1} \right)^\top.
\end{align}
As the set of monomials $\{x^k\}_{k=0}^{n-1}$ and the set of monic orthogonal polynomials $\{ p_k(x)\}_{k=0}^{n-1}$ form bases of the space of degree $n-1$ polynomials, there exist a change of bases matrix $L_{n}$ such that it is an $n \times n$ lower-triangular matrix with all diagonal entries $1$ and satisfies \( \bm{p}_n(x) = L_n \bm m_n(x) \).
We denote the $(n-m) \times (n-m)$ upper left block of $L_{n}$ by $L_{n}^{[n-m]}$ that satisfies \( \bm{p}_{1}(x) = L_n^{[n-m]} \bm m_{1}(x) \).

Next, we define the $n \times (n-m)$ moment matrix with respect to $\mu$,
\[
G^{(1)}[\mu] = \int_{\mathscr I} \bm m_n (x)\bm m_{1} (x)^\top\,\dd\mu(x)=[\mu_{i+k}]_{\substack{0\le i\le n-1\\ 0\le k\le n-m-1}}.
\]
Then, one can readily observe that 
\begin{equation} \label{app key eqn}
\int_{\mathscr I} \bm p_n (x) \bm p_{1} (x)^\top \,\dd\mu = L_n G^{(1)}[\mu] \left( L_n^{[n-m]} \right)^\top.
\end{equation}
However, orthogonality \eqref{eq:orthog} simply implies that
\[
\int_{\mathscr I} \bm p_n (x) \bm p_{1} (x)^\top \,\dd\mu = \begin{pmatrix}
    \bm{H}_{n-m}[\mu]\\
    \bm{0}
\end{pmatrix} \in \mathrm{Mat}_{n \times (n-m)} \R,
\]
where \(\bm{H}_{n-m}[\mu] = \mathrm{diag}(h_0, h_1, \cdots, h_{n-m-1})\) (in the main text, this was called $\bm H_1$, see \eqref{defn of H} and \eqref{defn of M, H 4 blocks}).
Thus, the mixed-moment matrix $\bm{M}$ for a vector \eqref{app bordered hankel vector} and measures \eqref{app bordered hankel measures} is given by
\begin{align} \label{app bordered hankel M}
\bm{M} &= \left(\int_{\mathscr I} \bm p_n(x) \bm p_{1}(x)^\top \dd\mu(x), \bm p_n(z_1), \cdots, \bm p_n (z_m) \right) = \begin{pmatrix}
\bm H_{n-m}[\mu] & \bm M_{12} \\[2pt]
\bm 0    & \bm M_{22}
\end{pmatrix},
\end{align}
where $\bm M_{12} = \left(\bm p_1(z_1)\ \cdots\ \bm p_1(z_m)\right)$ is an $(n-m) \times m$ matrix and $\bm M_{22} = \left(\bm p_2(z_1) \cdots \bm p_2(z_m)\right)$ is an $m \times m$ matrix.

Since the bordered Hankel matrix \eqref{borderd hankel defn} can be written as
\[
\bm H_n^{B}[\mu;z_1,\cdots,z_m]
= \begin{pmatrix} G^{(1)}[\mu] & \bm m_n (z_1) & \cdots & \bm m_n (z_m)\end{pmatrix},
\]
using \eqref{app key eqn} and the first equality of \eqref{app bordered hankel M}, one can readily show that
\begin{equation}
\bm{M} = L_n \bm H_n^{B}[\mu;z_1,\cdots,z_m]\
\begin{pmatrix}
L_n^{[n-m]} & \bm 0\\
\bm 0 & I_m
\end{pmatrix}^\top.
\end{equation}
Combining it with the second equality of \eqref{app bordered hankel M}, we obtain the desired identity,
\[
\det\bm H_n^{B}[\mu;z_1,\cdots,z_m]= \det \bm{H}_{n-m}[\mu] \det \bm M_{22}.
\]
\end{proof}

\end{document}
